% ============ 第二部分：第 4--8 章（技术核心，含内联证明链） ============

\section{思想脉络（二）：把爆破化为配平问题的第一性原理}

\subsection{本质问题}

第 3 章已经把困难的位置指了出来：能不能造出爆破解，取决于能不能让残差 $f=\partial_{t}u+(u\cdot\nabla)u-\nu\Delta u+\nabla p$ 保持光滑。把这个判断再往前推一步，就得到整个技术方案的第一性原理。

若随便写一个在 $t\uparrow1$ 时发散的速度场，几乎必然得到一个同样发散的外力：一个刚性旋转、角速度按 $(1-t)^{-1/2}$ 增长的玩具涡旋，其压力取 $\frac12 a(t)^{2}(x^{2}+y^{2})$ 恰好抵消离心项，但残差仍然按 $(1-t)^{-3/2}$ 发散。可见“发明一个奇异涡旋”是容易的，难的是让残差里的各个发散分量恰好相消。于是工作真正的形状被确定下来：不是解一个方程，而是设计一场精确的抵消。

\begin{keybox}
\textbf{第一性原理}\quad \NS{} 爆破构造的困难不在动力学而在代数：要在速度、压力、对流、黏性四项各自的量级都趋于无穷时，安排它们的和连同一切导数保持有界。全部技巧都是在为这场抵消寻找足够的自由度。
\end{keybox}

\subsection{仿射背景上的精确波}

抵消的自由度从哪来？来自非线性项本身的“不自作用”性质。考虑一段仿射背景 $u_{\text{old}}(x,t)=D(t)x$，其中 $\operatorname{tr}D=0$（保证无散），再叠加一列平面波
\begin{equation}
w=a(t)\,\phi\big(\xi(t)\cdot x\big),\qquad a\perp\xi,
\end{equation}
其中 $a\perp\xi$ 使 $\nabla\cdot w=(a\cdot\xi)\phi'=0$。此时波的对流项
\begin{equation}
(w\cdot\nabla)w=a(a\cdot\xi)\,\phi\,\phi'=0
\end{equation}
因为 $a\cdot\xi=0$ 而恒为零。这就是 Craik--Criminale 的精确波：波的二次自作用消失，波因而满足一个线性方程。非线性项里最麻烦的那一块被结构性地点掉了，剩下的动力学变成可控的线性系统。这类结构在流体中并不罕见——离心不稳定性、Rayleigh--Taylor 不稳定性都可以用精确波来承载。

\subsection{振幅方程：逐项推导}

为看清背景应变 $D$ 如何给波注入能量，把 $u=Dx+w$ 代入无外力 Euler 方程，并减去背景自身满足的方程。以下逐步列出存活项，这一步是后面一切定量估计的起点。

设背景满足 Euler：$\partial_{t}(Dx)+(Dx\cdot\nabla)(Dx)+\nabla p_{\text{old}}=0$。由于 $(Dx\cdot\nabla)(Dx)=D^{2}x$ 与 $\partial_{t}(Dx)=\dot Dx$ 都是关于 $x$ 线性的，背景压力必须也是 $x$ 的二次型，这与后面波的压力恰好相容。把 $u=Dx+w$ 代入并减去背景方程，剩余项为三部分：

其一，时间导数项 $\partial_{t}w=\dot a\,\phi+a\,\phi'\,(\dot\xi\cdot x)$。其二，波对流背景 $\big(w\cdot\nabla\big)(Dx)=Dw$，以及背景对流波 $\big(Dx\cdot\nabla\big)w=\phi'\,(Dx\cdot\xi)a$。其三，波的自身对流项已由 $a\cdot\xi=0$ 消去，只剩下压力项 $\nabla p_{w}=q(t)\,\xi\,\Phi(\xi\cdot x)$，其中 $\Phi'=\phi$，故压力梯度沿 $\xi$ 方向。

采用相位传输条件 $\dot\xi=-D^{\top}\xi$，则 $\partial_{t}w$ 中沿 $\phi'$ 的分量 $a\,(\dot\xi\cdot x)$ 与背景对流波项 $a\phi'\,(Dx\cdot\xi)$ 恰好相消——注意 $\dot\xi\cdot x=-(D^{\top}\xi)\cdot x=-\xi\cdot(Dx)$，两者之和为零。余下的 $\phi$ 系数项给出
\begin{equation}
\dot a=-Da-q\,\xi ,
\label{eq:coeffphi}
\end{equation}
其中 $q$ 是待定的压力标量。为使波的无散约束 $a\cdot\xi=0$ 在演化中保持，需要对它求时间导数并令其为零：$\dot a\cdot\xi+a\cdot\dot\xi=0$。代入 $\dot\xi=-D^{\top}\xi$，得 $a\cdot\dot\xi=-(Da)\cdot\xi$；再代入式 \eqref{eq:coeffphi}，得 $\dot a\cdot\xi=-(Da)\cdot\xi-q|\xi|^{2}$。两者相加为零，于是压力标量被唯一确定为
\begin{equation}
q=-2\,\frac{\xi\cdot Da}{|\xi|^{2}},
\end{equation}
而振幅满足
\begin{equation}
\dot a=-Da+2\,\frac{\xi(\xi\cdot Da)}{|\xi|^{2}}
=-\Big(I-\frac{2\,\xi\otimes\xi}{|\xi|^{2}}\Big)Da.
\label{eq:ampeq}
\end{equation}
矩阵 $I-\frac{2\xi\otimes\xi}{|\xi|^{2}}$ 是沿 $\xi^{\perp}$ 平面的反射：压力把 $-Da$ 中沿 $\xi$ 的分量抹去，又因 $\xi$ 自身在转动而反向补回。范数满足
\begin{equation}
\tfrac12\partial_{t}|\xi|^{2}=-\xi\cdot D\xi,\qquad \tfrac12\partial_{t}|a|^{2}=-a\cdot Da .
\end{equation}
这两式给出一个干净的判据：$\xi\cdot D\xi<0$ 使波数变大，但它本身并不保证振幅增长；振幅是否增长，本质上是背景 $D$ 的稳定性问题，可以从式 \eqref{eq:ampeq} 读出。这就是“精确波”这一想法的全部内容——把一个偏微分方程问题归约为一对常微分方程，再在这对常微分方程上寻找不稳定线。

\subsection{两种截然不同的用法}

这套归约在两篇工作里扮演了相反的角色，这个对照是理解 \NS{} 构造的关键。

在 Boussinesq 中，被放大的波本身就是奇点。取自由的 $D=0$、稳定的温度梯度 $G=-Ae_{2}$（$A>0$ 对应冷重流体压住暖轻流体的不稳定层结），写 $\zeta=r\,e(\varphi)$，则振幅方程给出二阶系统
\begin{equation}
\frac{d}{dt}\begin{pmatrix}\Theta\\ \Omega\end{pmatrix}
=\begin{pmatrix}0 & \dfrac{A\sin\varphi}{\lambda r}\\[2mm] \lambda r\sin\varphi & 0\end{pmatrix}
\begin{pmatrix}\Theta\\ \Omega\end{pmatrix},
\qquad\text{特征值}=\pm\sqrt{A}\,\sin\varphi .
\end{equation}
生长率 $\sqrt{A}\sin\varphi$ 与波数 $\lambda$ 无关，但原点处的梯度 $\nabla\vartheta(0,t)=\lambda\,\Theta\,\zeta$ 与 $\lambda$ 成正比。于是，一个振幅极小、频率极高的波可以在温度本身保持有界的同时，制造出无界的梯度——恰好是“$\|\theta\|_{\infty}$ 有界而 $\|\nabla\theta\|_{\infty}\to\infty$”所需的那种形状。代价是随之而来的剪切 $\Omega$，需要用共同旋转把它“转向”（旋转背景梯度与波矢，使实验室坐标下的 $\zeta_{1}$ 反号，把 $\Omega$ 拉回零），之后才能在其上叠加下一层更细的波。每一层因此经历“生长、过渡、转向、保持”四阶段，层间时间可求和。

在 \NS{} 中，奇点不再由波承担。构造者先写下一条光滑的自相似坍缩涡旋柱，让它自己去发散；被放大的波只是工具，任务是提供坍缩本身缺掉的那份外力。这个分工的落地方式就是下一节要讲的 Reynolds 应力。

\subsection{Boussinesq 级联的全景：从一层到无穷层}

上一小节给了单层的增长率，但“一层”不足以产生奇点，需要把无穷多层塞进有限时间。这一步的组织方式是：$\lambda_{q}=\lambda_{q-1}^{Q_{*}+q}$ 的超指数频率表、$\ell_{q}=\lambda_{q-1}^{-3}$ 的包络半径，以及 $\sigma_{q}\ge\sigma_{q-1}^{3}$ 的增长尺度。每层的生长时长约为 $\sigma_{q-1}^{-1}$，而 $\sigma_{q}$ 按立方增长，故各层时长求和收敛到一个有限的 $T_{*}$。于是无穷多层在 $T_{*}$ 之前落位，梯度在 $T_{*}$ 处发散。

要让这个级联真的收敛到“梯度发散而温度有界”，还需处理两件事。其一是增益不能被后续层抵消：所有梯度方向都落在以第一层方向为中心的一个锐角锥内，$\zeta_{q}$ 的方向被共同旋转“转向”后仍与原方向保持夹角下界，于是投影到 $-e_{1}$ 上的和仍有正的下界，$\|\nabla\theta(t)\|_{\infty}\ge cA_{n}(t)\to\infty$。其二是各层不能互相破坏：由于 $\ell_{q}\ll\ell_{q-1}$，更细的层被支撑在所有更粗层的仿射区域内，那里剧烈部分恰好是精确仿射的，高阶修正项为零，从而层与层之间的干扰被结构性排除。

这套“超指数频率＋锐角锥方向控制＋嵌套仿射区域”的三件套，是 Córdoba--Martínez-Zoroa 纲领的具体化身。理解了它，\NS{} 构造里那些看似琐碎的频率与阶数安排就有了出处——它们不是审美选择，而是收敛性的硬性要求。

\subsection{精确波为何是“引擎”}

把精确波称为引擎，不只是因为它提供增长，更因为它把“非线性爆破”这件难事，拆成了“线性稳定性”这件易事。爆破的时间、位置与速率，都可以从振幅 ODE 的特征值与相位方程直接读出：特征值 $\pm\sqrt{A}\sin\varphi$ 给出生长率，$\dot\xi=-D^{\top}\xi$ 给出波数随时间的走向，$\lambda\Theta\zeta$ 给出梯度的量级。换言之，一个原本要在偏微分方程层面讨论的问题，被压缩到几个标量 ODE 上。

这种压缩之所以可能，依赖两条结构性事实：波的极化与波矢正交（$a\perp\xi$，消去自作用），以及背景是仿射的（$\nabla(Dx)=D$ 是常数矩阵，使代数运算初等化）。一旦背景不再是仿射的，或波不再正交极化，压缩就会失效。于是这一工具的适用范围被划得很清楚：它服务于“局部仿射、局部平面波”的区域，其它区域要靠包络的截断与逐阶修正去处理。第 5、6 章分别把这两类区域用到核心与环域上。

\subsection{小结}

把这一章压缩成一句可记的话：精确波之所以有用，是因为它把“非线性”中最麻烦的自作用项结构性地消掉，从而把一类看起来无从下手的偏微分方程问题，变成一对可以手算稳定性的常微分方程；剩下的工作，就是决定让这对常微分方程的不稳定性去承担奇点，还是去承担外力。

\section{核心数学技术（一）：各向异性自相似坍缩涡旋}

\subsection{奇点的几何：一根收缩的细柱}

构造把奇点放在原点、时间归一化为 $t=1$。在柱坐标 $(r,\theta,z)$ 下，主导流 $u^{(0)}$ 是轴对称的，其分量只依赖 $(r,z,t)$。物理图像很直观：流体沿径向向内螺旋、同时沿轴向在 $z\approx0$ 分层面两侧向外排出；压力向轴心递减提供向心的主导力。角动量在向内输送中被黏性向外搬运，二者平衡决定了特征转速。

关键是自相似的标度。记 $\tau=1-t$，手稿取
\begin{equation}
\ell_{r}\asymp\tau^{1/2},\qquad \ell_{z}\asymp\tau^{1/2-h},\qquad 0<h<\tfrac1{100},
\label{eq:scales}
\end{equation}
其中 $\asymp$ 表示被与 $\tau$ 无关的正常数上下夹住。两个长度都趋于零，但
\begin{equation}
\frac{\ell_{r}}{\ell_{z}}\asymp\tau^{h}\to0,
\end{equation}
所以核心越来越细长，其体积为 $\ell_{r}^{2}\ell_{z}\asymp\tau^{3/2-h}$。

\subsection{速度标度与能量：局部无界与全局有界的相容}

与式 \eqref{eq:scales} 配套的速度量级为
\begin{equation}
|u_{\theta}^{(0)}|,\,|u_{z}^{(0)}|\asymp\tau^{-1/2-h},\qquad |u_{r}^{(0)}|=O(\tau^{-1/2}).
\end{equation}
现在直接估算核心的动能，看它为什么不发散：
\begin{equation}
E_{\text{core}}\;\sim\;\text{体积}\times(\text{特征速度})^{2}
\;\asymp\;\tau^{3/2-h}\cdot\tau^{-1-2h}
\;=\;\tau^{1/2-3h}\;\xrightarrow[\tau\to0]{}\;0 .
\label{eq:energy}
\end{equation}
因为 $h<1/100$，指数 $1/2-3h>0$ 仍为正，故核心动能趋于零——尽管速度在增长。这就是“局部速度无界、全局动能有界”之所以相容的算术根源：速度按 $\tau^{-1/2-h}$ 发散，但承载它的体积按 $\tau^{3/2-h}$ 收缩，二者相乘后能量反而变小。

\begin{keybox}
\textbf{一个可记的估算}\quad 无限速度之所以不违反有限能量，是因为它被压缩进了越来越小的体积。式 \eqref{eq:energy} 的指数 $1/2-3h$ 就是“能量不掉队”的全部秘密，$h$ 的上界 $1/100$ 正是为了保住这个不等式。
\end{keybox}

\subsection{两套 Reynolds 数}

用核心半径 $\ell_{r}$ 作长度、以角向与径向速度分别作速度，得
\begin{equation}
\mathrm{Re}_{\theta}=\frac{|u_{\theta}|\ell_{r}}{\nu}\asymp\tau^{-h}\to\infty,
\qquad
\mathrm{Re}_{r}=\frac{|u_{r}|\ell_{r}}{\nu}=O(1).
\end{equation}
角向 Reynolds 数无界增长，意味着流体在径向扩散时间内转了越来越多的圈；径向 Reynolds 数保持有界，意味着黏性始终与径向内流相抗衡。再看各项的变化率：
\begin{equation}
\frac{|u_{r}|}{\ell_{r}}=O(\tau^{-1}),\quad
\frac{|u_{z}|}{\ell_{z}}\asymp\tau^{-1},\quad
\frac{\nu}{\ell_{r}^{2}}\asymp\tau^{-1},
\qquad
\frac{\nu/\ell_{z}^{2}}{\nu/\ell_{r}^{2}}=\frac{\ell_{r}^{2}}{\ell_{z}^{2}}\asymp\tau^{2h}\to0 .
\end{equation}
前三个变化率都是 $\tau^{-1}$ 量级，说明径向扩散、径向对流与轴向对流同处主导平衡；轴向扩散相对径向扩散按 $\tau^{2h}$ 衰减而退居次席。这套尺度不是拍脑袋选的，而是“让主导项同阶、让次主导项可压”这一要求反解出来的。

\subsection{为什么必须打破反射对称}

若剖面关于 $z=0$ 严格对称，则角动量向中层面的轴向输运在 $z=0$ 处消失，且对称性迫使 $u_{z}(r,0,t)=0$，于是中层面附近既无转动放大、也无轴向速度的径向剪切可以补足。构造因此选取略不对称的轴向剖面：带一点向上的偏置、在 $z=0$ 处留一个非零的 $u_{z}$。这使“转动放大较弱的层”与“轴向剪切消失的层”在位置上分开，由 $u_{z}$ 的径向剪切在中层面附近补足所需的放大；稍远处，转动机制本身就够用。这一处看似琐碎的对称性破坏，是核心侧壁能把剪切喂给外侧脉冲的前提。

\subsection{与刚性定理假设的对应}

后文第 9 章要讨论的 Constantin--Ignatova--Vicol 刚性定理\cite{civ}，其假设正好抓住了本节的两个特征：其一，速度的角向平均 $\mathcal{P}u$ 满足各向异性的 Type II 界，即径向分量被 $(1-t)^{-1/2}$ 控制、轴向与角向分量被 $(1-t)^{-1/2-h}$ 控制，且每求一次 $r$ 导数多付 $\tau^{-1/2}$、每求一次 $z$ 导数只多付 $\tau^{-1/2+h}$；其二，速度在一个半径为 $c\sqrt{\tau}$ 的坍缩核心区严格轴对称。这两条与式 \eqref{eq:scales} 及核心构造一一对应。把这个对应先记住，第 9 章的结论才读得顺。

\subsection{核心内部的剖面方程与压力平衡}

标度回答的是“多快、多小”，剖面方程回答的是“什么形状”。在核心最内侧，主导动量平衡由一组关于 $(r,z)$ 的定常型方程给出；把它写成自相似变量后，径向与轴向的平衡各自成为一个常微分方程的边值问题（附录 B 专门处理轴附近的解析剖面与延拓）。物理上，压力向轴心递减，提供向内螺旋所需的向心力；轴向压力的梯度在中层面附近指向 $z=0$，与轴向输运和黏性共同维持轴向排出。由于不可压，向内汇集的流体必须沿轴向排出，否则会在轴心堆积——正是这条约束把“向内”、“旋转加速”、“轴向排出”三者锁在一起：径向内流携带角动量靠近轴心，转动因而加快；无扭矩时角动量 $ru_{\theta}$ 守恒，而实际核心中黏性把角动量向外搬走，特征转速由“向内搬运”与“黏性流失”的平衡决定。

这套剖面方程的存在性本身依赖附录 B 的解析性论证：只有在轴附近取到实解析的剖面，内层才能与外部流光滑对接。这也是“解析性”在本构造中第一次以正面身份出现——第 9 章刚性定理里，外力一旦在奇点附近解析就会逼出正则，而这里解析性只用于核心内部的剖面，两者作用方向相反，不可混为一谈。

\section{核心数学技术（二）：振荡脉冲与平均动量通量}

\subsection{环域里的“缺力”}

把自相似核心与一个速度随半径递减的光滑外部流对接，会在两者之间的环域留下一个不平衡。单看背景流，这个残差在 $t\uparrow1$ 时发散，无法充当所需的光滑外力。于是必须让流体自身的运动补上这份缺掉的动量输运——具体做法，是在环域布置一族空间振荡的脉冲。

\subsection{零角平均脉冲如何产生内部力}

脉冲的关键性质是：它们的角向平均为零，但其二次自作用不为零。对速度分量做角向平均，记 $\langle\cdot\rangle=\frac{1}{2\pi}\int_{0}^{2\pi}\cdot\,d\theta$，则对流项的平均包含
\begin{equation}
\big\langle (w\cdot\nabla)w\big\rangle
=\big(\langle w\rangle\cdot\nabla\big)\langle w\rangle
+\big\langle (w-\langle w\rangle)\cdot\nabla\,(w-\langle w\rangle)\big\rangle .
\end{equation}
第一项因 $\langle w\rangle=0$ 而消失，第二项一般不为零：两个平均为零的振荡分量，其乘积的平均不必为零。这一项正是 Reynolds 应力 $\langle w\otimes w\rangle$ 的散度，它充当作用在背景流上的内部力。选择脉冲的排布，使其提供的 Reynolds 应力恰好抵消背景残差的奇异部分，剩下的余项就能光滑地穿过 $t=1$。这一步是本项工作区别于前人的核心创新。

\begin{keybox}
\textbf{核心创新的一句话版}\quad 用平均为零的振荡脉冲，靠它们乘积的平均（Reynolds 应力）去补足坍缩涡旋在奇点处缺掉的那份动量输运；脉冲自己不发散，却替背景流承担了发散的外力。
\end{keybox}

\subsection{脉冲何以增长：离心不稳定性}

脉冲需要从“指数小”逐渐生长到足以承担通量。其生长机制是经典的离心不稳定性，可以追溯到 Leibovich--Stewartson 与 Billant--Gallaire 的离心不稳定判据\cite{leibovich,billant}。在环域中，角速度随半径向外递减，于是被径向推移的流体元处在不利于稳定的角动量分配中，扰动被进一步推离原位，形成指数增长。第 4 章的精确波归约在此再次出场：把脉冲看作叠加在缓慢变化的仿射背景上的精确波，其振幅 ODE 给出增长指数；脉冲的振荡性保证它对背景的“平均”效应正是我们要的 Reynolds 应力，而其局部起伏不破坏背景的光滑性。

\subsection{黏性阻尼与激活窗口}

与 Euler 方程不同，\NS{} 的脉冲要额外背负一项黏性阻尼。对正弦载波的精确波，方程中会多出一项 $-\nu|\xi|^{2}a$：频率越高，阻尼越强。这意味着脉冲不能无限早地布置下去——若种下得太早，它会在生长到有用规模之前就被黏性抹平；若种得太晚，则来不及长到需要的振幅。因此脉冲必须有一个精心安排的激活时刻，并由一个指数小的外力种子启动（手稿提到“指数小的外力种下每个脉冲”）。种子的量级必须与允许的生长时间相匹配：生长率与所需增益给定时间；时间给定允许的种子大小。

这解释了一个乍看奇怪的设计：脉冲要靠外力启动，而这个外力是构造本身安排的一环，且小到不影响外力整体的光滑性。在“残差博弈”的账本上，这是一笔划算的交易——用一个指数小的种子换来一个在奇点附近承担全部发散的正则化机制。

\subsection{离心不稳定性的物理判据}

离心不稳定性有明确的判据形式，可追溯到 Rayleigh 判据及其推广：在轴对称旋转流中，若角动量的径向分布向外递减得足够快（比 $r^{-2}$ 更快），则流是离心不稳定的。Leibovich--Stewartson 与 Billant--Gallaire 的准则把这套判据推广到非轴对称扰动。手稿安排核心外侧的角速度分布，使其恰好落在不稳定一侧；于是脉冲一旦被种下，其振幅增长由该判据保证的指数率驱动。

有一点值得注意：这里的稳定性判据与第 4 章的精确波振幅方程是同一件事的两种表述。Rayleigh 判据看的是角动量分布，振幅方程看的是 $D$ 的特征值——两者给出同一个“是否增长”的答案，只是语言不同。这种“同一事实、两种语言”的现象在本领域很常见，也是它难以初读的原因之一。

\subsection{尺度匹配}

脉冲的尺度必须与核心的坍缩尺度相配。核心以 $\ell_{r}\asymp\tau^{1/2}$、$\ell_{z}\asymp\tau^{1/2-h}$ 收缩，环域因此也在收缩；脉冲若要持续供给恰好抵消奇异残差的通量，其波长、振幅与激活时刻都必须随 $\tau$ 同步安排。手稿的作法是让脉冲“分族”布置——在固定高度上按整环排布，在径向—轴向上按局部小片排布——使得两类扰动分别承担径向与轴向的动量通量，且在给定的角向平均下相互补充。把这几套尺度对齐，是第 7 章逐阶修正能够启动的前提。

\subsection{与 Córdoba--Martínez-Zoroa 的差别}

同样是“被放大的扰动”，两处角色不同。Córdoba--Martínez-Zoroa 的放大扰动（涡层、振荡层）本身把解推向更奇异，奇点是放大的累积效果；OpenAI 的脉冲则服务于一个预先写定的光滑坍缩，其角色是把发散移交给 Reynolds 应力，从而让外力保持光滑。这个差别解释了为什么 \NS{} 的构造需要处理“脉冲—背景—外部流”三段式，而 IPM/Boussinesq 的构造更像是一串首尾相接的层。

\subsection{脉冲的几何：整环与局部小片}

手稿的脉冲不是各向同性的随机扰动，而是有明确几何。在固定高度上，脉冲包络铺满完整的环（其径向与角向扰动速度在环内变化），构成“整环族”；在径向—轴向截面上，脉冲则聚成一系列局部小片，每片是某个整环的横截面，另构成“小片族”。两族扰动的极化方向不同（一族给径向、一族给轴向），因此在给定的角向平均下能够互补：把它们的 Reynolds 应力 $\langle w\otimes w\rangle$ 加起来，正好覆盖背景残差需要被抵消的那几个分量。

这套几何安排有两个好处。其一是“平均为零、二次不零”这一性质得到保证：整环在小片尺度上仍然振荡，故其自作用的平均不为零，而它对背景的线性影响因平均为零而消失，从而只留下我们要的通量项。其二是与轴向外流的兼容：小片沿轴向梯度较大，可在轴向动量平衡中扮演主导，而整环在角向的变化则服务于角向平衡。附录 C 的“容许应力锥”就是把这套几何翻译成一个可以被逐点满足的应力实现问题——给定想实现的应力张量，问是否存在一族这样的脉冲去实现它，答案是肯定的（在一个锥内）。

\section{核心数学技术（三）：逐阶修正与外力光滑性}

\subsection{残差的阶次结构}

单个 ansatz 不可能让残差严格为零；它总会留下误差。这里的处置办法与第 3 章的纲领一致：把残差按阶次组织，用一个无穷序列的修正把误差不断推到更高阶，使其在每一个导数意义上都趋于零。手稿的第三至第十章就是这条修正链的展开，其骨架可以分成四段：构造主导流（第四章）、把基流修正到任意阶（第五章）、用辅助环面把振荡实现出来并修正残余应力（第六、七章）、以紧支撑的均值修正与残差改进收尾并推广到全空间（第八至十章）。

\subsection{与 Alpöge--Buckmaster 的机制对照}

Alpöge--Buckmaster 的 Boussinesq 构造里，有一个可以直接类比的机制：在包络变化之处，波 ansatz 会留下振荡误差，这些误差被显式修正到阶 $J_{q}=2k_{q}+8$，而 $k_{q}\to\infty$。让修正阶数随层号趋于无穷，是 Córdoba--Martínez-Zoroa 的“阶数随尺度上升”原则的延续，也正是外力级数能逐导数收敛的原因。\NS{} 的逐阶修正遵循同一逻辑，只是被修正的对象更复杂——除了背景残差，还包括脉冲之间、脉冲与背景之间的交叉项。

\subsection{构造链条与三则附录}

手稿把技术细节分摊到三个附录：附录 A 做径向矩匹配并构造热核型的“外部”，使核心能与一个速降的外部流对接；附录 B 给出轴附近的解析剖面及其延拓，为核心内层的正则性服务；附录 C 处理容许应力锥的实现，即把“想实现的 Reynolds 应力”翻译成“能实际布置出来的脉冲族”。这三块合起来，才使前面几章的标度与通量论证落地为一个真正全局定义的光滑构造。

\subsection{从主导流到全局构造：四步链条}

手稿第三至第十章的展开，可以归纳为一条四步链条，每一步解决一个具体障碍。

第一步（第四章）是先造出“在核心内精确成立、在外部光滑”的主导流：核心给出自相似坍缩，外部给出速降的缓流，两者在环域对接。这一步留下一个问题——环域的残差发散。

第二步（第五章）把基流修正到任意阶：以逐阶展开的方式，把残差里可消去的部分尽量在 ansatz 层面消去，使剩余误差的阶数尽量高。这一步的作用是降低后续修正的负担，但它无法处理环域中真正发散的那一块。

第三步（第六、七章）用辅助环面把振荡“实现”出来：先在一个模型环面 $\mathbb{T}^{2}$ 上把脉冲的构造与估计做干净，再把它搬回柱坐标下的环域，并用显式修正消除 ansatz 留下的振荡误差。引理 8.6 正是这一阶段反演方向导数算子的工具。

第四步（第八、九、十章）收尾：用紧支撑的均值修正把全局支撑压紧，改进残差的余项，最终把构造推广到整个空间并得到紧支外力。

这条链条的每一环都以“残余误差的阶数”为接口：前一步的输出是后一步的输入，接口上的阶数必须单调上升，这正是“修正阶数随层号趋于无穷”在本构造中的对应形态。读懂“阶数在链条上如何流动”，是读懂全篇的关键。

\subsection{外力为何光滑却不解析}

有一处细节必须在这里点明，因为它直接连到第 9 章的刚性定理。手稿的外力被定义为推进残差，它 $C^{\infty}$ 光滑，但并不是实解析的：一方面，它在空间上紧支撑，紧支撑的非零实解析函数不存在，故它在全空间上不可能解析；另一方面，它在奇点邻域也非解析，因为它由紧支撑的截断构造拼出。这个“光滑但不解析”的性质当时看起来只是技术副产品，事后却被证明是构造得以成立的必要条件之一——第 9 章会看到，若外力在奇点附近实解析，构造将被迫正则，与爆破相矛盾。

\section{核心数学技术（四）与验证争议：引理 8.6 的形式化落差}

\subsection{引理 8.6 的位置与作用}

引理 8.6 位于第八章“紧支撑均值修正”。它的作用是把一个局部问题降到环面 $\mathbb{T}^{2}=\mathbb{R}^{2}/\mathbb{Z}^{2}$ 上，并在那里反演一个方向导数算子。设 $F$ 是参数 $p$（外部空间变量）与环面变量 $y$ 的函数，固定方向 $v_{t}$ 定义算子 $N=v_{t}\cdot\nabla_{y}$。要在修正中“除以 $N$”解出某个势，就需要控制 $N^{-1}$ 对 $F$ 各阶导数的放大倍数。引理 8.6 给出的正是这样一个反演估计。它是整条修正链里一个不起眼却必须成立的环面估计。

\subsection{内联证明链：除子界、反演乘子与可和性}

在手稿的语境里，这个估计的证明很短，且每一步都可手推，值得完整走一遍，因为它把“数论里的除子界”和“二维 Fourier 级数的可和性”直接焊在了一起。

先看 Fourier 侧。把 $F$ 在环面 $\mathbb{T}^{2}$ 上展开，$\widehat{F}(k)$ 记频率 $k\in\mathbb{Z}^{2}$ 处的系数。算子 $N=v_{t}\cdot\nabla_{y}$ 在频率 $k$ 上作用为乘上 $i\,(v_{t}\cdot k)$，故其形式逆 $N^{-1}$ 的象征是 $\big(i\,(v_{t}\cdot k)\big)^{-1}$，放大倍数即
\begin{equation}
\big|v_{t}\cdot k\big|^{-1}.
\end{equation}
这是估计的第一步：只要方向 $v_{t}$ 固定且非零，频率 $k$ 离方向 $v_{t}$ 的垂直平面越近，$|v_{t}\cdot k|$ 就越小，放大就越厉害。手稿在此调用一个除子型的界，把这一放大用频率的总大小兜住：
\begin{equation}
\big|v_{t}\cdot k\big|^{-1}\;\le\;C\,(1+|k|).
\label{eq:divisor}
\end{equation}
式 \eqref{eq:divisor} 的含义是：最坏情形（$k$ 几乎垂直于 $v_{t}$）下，放大倍数被频率本身线性控制——因为整数格点靠近某个过原点的直线时，偏离量有其算术下界。

第二步是可和性。求导在频率侧对应乘上 $|k|$。设要求控制输出到 $m$ 阶导数，而输入比输出多付 $D$ 阶导数，则每个频率分量的贡献为
\begin{equation}
|k|^{m}\,\big|v_{t}\cdot k\big|^{-1}\,|k|^{-(m+D)}
\;\le\;C\,|k|^{m}(1+|k|)\,|k|^{-(m+D)}
\;\lesssim\;(1+|k|)^{1-D},
\label{eq:summ}
\end{equation}
其中第二步用了除子界 \eqref{eq:divisor}。于是问题归结为：$D$ 取多少，才能让 $(1+|k|)^{1-D}$ 在二维格点上可和。二维的判据是
\begin{equation}
\sum_{k\in\mathbb{Z}^{2}}(1+|k|)^{-s}<+\infty\iff s>2,
\end{equation}
因为半径 $|k|\sim R$ 的格点数按 $R$ 增长，与 $(1+|k|)^{-s}$ 相乘得到 $\sim R^{1-s}$，而 $\sum_{R}R^{1-s}$ 收敛当且仅当 $1-s<-1$，即 $s>2$。记衰减式 $(1+|k|)^{1-D}=(1+|k|)^{-s}$ 中的 $s=D-1$，则二维可和要求 $s>2$，即 $D>3$；整数下最小可行值即为 $D=4$，此时衰减恰为 $(1+|k|)^{-3}$，可和，但距不可和的 $(1+|k|)^{-2}$ 只差一阶、没有任何余量。这就是“比输出多四个输入导数”的来源。若换到三维环面，判据变为 $s>3$，即 $D>4$，最小整数为 $D=5$——所以这条引理必须在二维环面上陈述，多一维就不够。

至此可和性成立，进而给出估计 \eqref{eq:summ} 对应的算子界。手稿随后补一句：该乘子与一切“慢变量”系数导数可交换，且保持函数的实性；再补一句：固定零频系数后解唯一。这三句话合起来，就是引理 8.6 的完整陈述。

\begin{remark}
这套证明有一个容易被忽略的结构特征：唯一真正用到“数论”的地方是除子界 \eqref{eq:divisor}，其余部分只是 Fourier 级数的初等记账。之所以要绕道频率侧，是因为方向导数算子 $N=v_{t}\cdot\nabla_{y}$ 在物理侧不可逆（沿 $v_{t}$ 方向为常数、沿垂直方向属零空间），而在频率侧它的逆只是一个逐频相乘的因子；代价是这个因子在 $v_{t}\cdot k$ 接近零时爆炸，需要除子界来控制。这类“在频率侧求逆、用算术界兜底”的手法，属于拟微分算子记帐的常见套路，但它在这里被安排在一条纯初等的证明里，用起来很轻。

还须补上关于这条界本身的一点：它不是自动成立的记账琐事，而是一条带算术内容的输入。对有理方向，$|v_{t}\cdot k|^{-1}$ 一致有界，远强于所需；对无理方向，Dirichlet 逼近定理保证 $|v_{t}\cdot k|\lesssim|k|^{-1}$ 对无穷多个 $k$ 成立，最坏情形下增益因子 $(1+|k|)$ 被用满——这正是 $D=3$ 不足、$D\ge4$ 的原因；若方向具有更极端的有理逼近（Liouville 型），这条界本身会失效，任何固定导数额度都不足以收尾。因此 $m+4$ 不能被视为记账的自动推论：它依赖方向 $v_{t}$ 的算术性质，而该性质在构造里由对方向的选择决定。
\end{remark}

\begin{keybox}
\textbf{证明链的三步}\quad （1）逆算子的象征是 $|v_{t}\cdot k|^{-1}$；（2）除子界把它换成 $C(1+|k|)$；（3）输入比输出多给 $D$ 阶导数后，衰减为 $(1+|k|)^{1-D}$。二维格点的可和判据是 $s>2$，故 $D=4$ 时给 $(1+|k|)^{-3}$（$s=3$，恰可和、无余量），这正是“多四阶”的来源；三维判据是 $s>3$，需 $D\ge5$，所以引理只能在二维环面上陈述。
\end{keybox}

\subsection{落差：多四个导数还是多五个}

争议出现在自然语言证明与 Lean 形式化之间。手稿正文在引理 8.6 的开头给出如下论证：由除子界 $|v_{t}\cdot k|^{-1}\le C(1+|k|)$，对输入 $F$ 求的导数比所求输出多四个，即留下二维可和的 $(1+|k|)^{-3}$ 界，式 (8.19) 由此得证。但 Bastounis、Circelli 与 Hansen\cite{bch}指出，Lean 仓库中与式 (8.19) 最接近的定理 \texttt{norm\_derivativeWord\_inverse\_le}，其假设用的是“多五个导数”，即控制同一批环面输出导数时，Lean 版本要求输入多出 $m+5$ 阶而非 $m+4$ 阶。就命题强度而言，多要求一阶输入正则性是更强的假设，因此 Lean 该定理所证明的命题弱于手稿声称的命题；这正是研究者所称“自然语言论文声明了一个更强的反演估计”的所指。

联系上一小节的记账式 \eqref{eq:summ} 可以把这一差别算准：该式对一般的输入导数额度 $D$ 给出衰减 $(1+|k|)^{1-D}$，与 $D$ 取四或五无关地成立。代入 $D=4$ 得 $(1+|k|)^{-3}$（$s=3>2$，可和但无余量），代入 $D=5$ 得 $(1+|k|)^{-4}$（$s=4>2$，可和且有余量）；两者在二维环面上都收敛；不可和的只有 $D\le3$，其中 $D=3$ 恰为临界情形（衰减 $(1+|k|)^{-2}$）。因此落差之处不在可和性，而在命题的强弱：$m+4$ 是更强的估计，它用较少的输入正则性控制同一批输出导数；$m+5$ 则较弱。手稿第 96 页声称的 (8.19) 属于前者，而 Lean 中对应的定理证明的是后者。按 Bastounis 等人的表述，形式化在这里并未覆盖手稿声称的那条更强的估计。

以上落差此前只见于 Bastounis 等人的转述。本文进一步核对了公开仓库的一手源码，结论与转述一致。其一，仓库 \texttt{openai/}\allowbreak\texttt{NavierStokesAndEuler} 在 commit \texttt{f9e8bc5} 处的完整文件树共 2673 个条目，其中 \texttt{.lean} 文件 2659 个；在该仓库内检索声明名 \texttt{norm\_derivativeWord\_}\allowbreak\texttt{inverse\_le}，命中文件唯一，即 \texttt{NavierStokes/}\allowbreak\texttt{SmoothFamilyTorusInverse.lean}。其二，该文件共 1140 行、57334 字节，blob 哈希为 \texttt{c607eb6f2460}\allowbreak\texttt{a3ba55dc2218}\allowbreak\texttt{072cd504bf62}\allowbreak\texttt{53bf}，被质疑的定理声明位于第 1059 行：其两个假设项分别对 $F$ 的切片与交换切片要求 \texttt{xJet (w.length + 5)} 有界，即输入比输出多五个导数；结论侧的权重因子写作 $(\texttt{weight}\ k^{4})^{-1}$，指数为 4。其三，向下游看，同一文件第 1082--1136 行的 \texttt{mixedJet\_}\allowbreak\texttt{inverse\_bound} 在该段内共十处沿用 \texttt{w.length + 5}，并在第 1135 行调用上述声明；可见 $m+5$ 的口径并非孤立一处，而是贯穿了这条估计链。核验所得的代码与 Bastounis 等人所引代码块逐字一致（对方仅另加了红色标注）。因此落差的 Lean 一侧——形式化定理确按 $m+5$ 表述——是经一手源码确认的事实，而非仅凭二手转述；正文一侧的 $m+4$ 记述则依据 Bastounis 等人对手稿原文的转引。

由此产生的是一个具体的问题清单：其一，(8.19) 在其字面强度（$m+4$）下是否成立；其二，若它只以 $m+5$ 成立，那么修正链中按 $m+4$ 的强度调用它的位置，其输入函数是否恰好备有多余一阶的正则性可供支取；其三，Lean 定理的结论侧是否与 (8.19) 逐字对应——Bastounis 等人的表述是“控制同一批环面输出导数”，但输出界中常数的具体形式是否与正文一致，则仍需逐式对照原文确认（代码一侧的假设强度与权重指数已由本文核验，见上段）。需要说明的是，这处落差表明的是“形式化覆盖的命题弱于正文声称的命题”，并不等同于正文的 (8.19) 有误：按第 8.2 节复原的四阶记账，$m+4$ 的论证在简化模型中是自洽的，落差的实际后果取决于调用点处是否留有余量。三问都需要逐处核对正文与该引理的调用点，本文无法依据公开材料给出结论。

\subsection{SCI 谱系论证：为何“翻译”可能比停机问题更难}

Bastounis--Circelli--Hansen 的更大主张不止于这一个引理。他们区分两种自动形式化：第一种是把一份“自然语言证明”翻译成 Lean 并让 Lean 通过编译；第二种是把一份数学文本“语义忠实地”翻译成 Lean。他们指出，第一种的成功不能推出原自然语言证明正确，因为 AI 可以在翻译时悄然改换语义——或者给错误证明配上一个正确证明，或者把正确证明换成另一个正确证明，或者像本案这样，把“更强的声明”翻译成“更弱的定理”而仍然编过。

他们把语义忠实翻译的困难度放在可解性复杂度层级（Solvability Complexity Index，SCI）上：要消除自然语言数学文本的歧义（这是语义忠实翻译的必要条件），其难度在 SCI 谱系中高达无穷（$\mathrm{SCI}=\infty$），用他们的话说，比包括停机问题（$\mathrm{SCI}=1$）在内的任何计算问题都更难。由此得出的操作性结论是：OpenAI 的自然语言证明与其他被自动形式化的证明，不应仅凭 Lean 证书就获得信任，而应经受与其他证明同样的同行评审与审视。

评估这道落差时需要同时握住两侧的事实。按 OpenAI 公开的 Lean 仓库及第三方对仓库的独立计数，形式化包含六十余万行 Lean 代码、三万余条定理与引理，未留未证明的 \texttt{sorry} 占位，也未引入标准三公理之外的公理；但仓库自述的评审状态是“self-assessed（自评）”。Bastounis、Circelli 与 Hansen 强调，自评状态与独立验证是两回事，而引理 8.6 的样本说明：即便 Lean 全部编译通过，正文论证的某些步骤仍可能未被覆盖\cite{bch}。两方援引的都是公开可查的材料，分歧的解决有赖于逐条核对正文与形式化的对应关系；本文对此不作裁断。

\subsection{这一落差是否影响主定理}

诚实的回答是：目前无法从公开材料确定。需要区分两个层次的问题。其一是“引理 8.6 按其字面（$m+4$）是否成立”；其二是“即便该引理应按 $m+5$ 陈述，主定理的证明链是否会因此断裂”。第二种情形的可能性是存在的：一个中间引理的常数或导数阶数，未必是主定理链条上不可替换的一环——若其他步骤留有冗余，或该引理可以换用别的估计，主定理仍可能成立。反之，若该引理恰是拼接处唯一的支点，落差就会传导到主定理。这个判断需要把引理 8.6 在整条修正链中的调用点逐一梳理，超出本文范围，留作开放问题（见第 11 章）。

从方法论上说，这件事给出的教训比它给出的结论更重要。自然语言证明的检验，历来靠“读者复算”；形式化证明的检验，靠内核编译。当两者之间存在一层 AI 自动翻译时，“复算”与“编译”都不再直接作用于原文——这就是 Bastounis 等人所说的断层。对本领域而言，一个可行的折中是把形式化的对象从“全文”缩小到“关键节点”：对每个支撑主定理的核心引理，独立写一份与正文严格对应的形式化陈述并加验收，而不是依赖一次性把整篇论文译成代码。第 11 章的开放问题会把这一条列为具体建议。
